\section{Full System Implementation}

This chapter describes the complete implementation of the pose estimation system, spanning the Python-based backend pipeline and the native \texttt{visionOS} application for Apple Vision Pro. The system establishes a client-server architecture where the Vision Pro captures RGBD frames and transmits them to a Python backend that performs ArUco marker detection, monocular depth estimation or RealSense depth acquisition, ROI masking, and 6D pose estimation with head tracking correction.

\subsection{System Architecture Overview}

The complete system consists of three primary subsystems:

\begin{enumerate}
    \item \textbf{VisionOS Client} -- Runs on Apple Vision Pro, captures frames, manages user interaction via gaze and pinch gestures, and renders 3D pose overlays in AR
    \item \textbf{Python Backend Pipeline} -- Processes incoming frames through a unified computer vision pipeline supporting both hardware (RealSense) and AI-based (Transformers) depth estimation
    \item \textbf{Pose Estimation API} -- External service providing 6D object pose from RGBD inputs and CAD models
\end{enumerate}

\begin{figure}[H]
    \centering
    \fbox{\parbox{0.9\linewidth}{System Architecture Diagram Placeholder}}
    \caption{Overview of the complete system architecture showing dataflow between Vision Pro, Python backend, and pose estimation service.}
    \label{fig:system_arch}
\end{figure}

The dataflow operates as follows:

\begin{enumerate}
    \item Vision Pro captures RGB frames via main camera
    \item User selects region-of-interest (ROI) using gaze + pinch gesture
    \item Client retrieves camera intrinsics from backend or uses device intrinsics
    \item Client sends cropped RGB frame and intrinsics to backend via HTTP
    \item Backend processes frame through CV pipeline (ArUco detection, depth estimation, masking)
    \item Backend forwards RGBD + intrinsics + CAD model to pose estimation API
    \item Pose API returns 4×4 transformation matrix $\mathbf{T}_{\text{object}}^{\text{camera}}$
    \item Backend applies head pose correction: $\mathbf{T}_{\text{corrected}} = \mathbf{T}_{\text{head}} \cdot \mathbf{T}_{\text{object}}$
    \item Client receives pose, converts to RealityKit coordinates, and renders CAD overlay
\end{enumerate}

\subsection{Python Backend Pipeline}

\subsubsection{Unified Computer Vision Module}

The core processing pipeline is implemented in \texttt{computer\_vision\_pipeline.py}, a GPU-accelerated module supporting dual-mode depth acquisition. This design allows seamless switching between hardware depth sensors (Intel RealSense D435/D455) and AI-based monocular depth estimation (Depth-Anything-V2 transformer model).

\textbf{Key Features:}
\begin{itemize}
    \item \textbf{ArUco Marker Detection} -- Detects 4×4 fiducial marker grid (DICT\_4X4\_50)
    \item \textbf{Pose Estimation} -- Solves PnP using all detected markers with IPPE refinement
    \item \textbf{Dual Depth Support} -- RealSense hardware or Transformers inference
    \item \textbf{ROI Mask Extraction} -- HSV color-based segmentation with configurable tolerances
    \item \textbf{GPU Acceleration} -- CUDA support for depth inference (50-100ms on GPU vs 2-5s on CPU)
    \item \textbf{Time Synchronization} -- All outputs timestamped for latency analysis
    \item \textbf{PNG Encoding} -- Lossless depth encoding for API compatibility
\end{itemize}

\subsubsection{Depth Acquisition Modes}

\paragraph{Mode 1: Transformers-based Monocular Depth Estimation}

When \texttt{use\_realsense=False}, the pipeline employs the Depth-Anything-V2 model\cite{depth_anything_v2}, a state-of-the-art transformer architecture for relative depth prediction from single RGB images.

\textbf{Processing Pipeline:}
\begin{enumerate}
    \item Convert BGR frame to RGB PIL image
    \item Run inference through HuggingFace pipeline (GPU accelerated)
    \item Normalize output depth map to 0-255 range: $D_{\text{norm}} = 255 \cdot \frac{D - D_{\min}}{D_{\max} - D_{\min}}$
    \item Generate disparity via inversion: $\text{Disp} = 255 - D_{\text{norm}}$
    \item Encode as PNG base64 for transmission
\end{enumerate}

\textbf{Performance:} 50-100ms per frame on CUDA-enabled GPU (RTX 3060/4060), 2-5s on CPU

\paragraph{Mode 2: RealSense Hardware Depth}

When \texttt{use\_realsense=True} and RealSense camera is available, the pipeline uses metric hardware depth via the \texttt{RealSenseToAVPAligner} adapter.

\textbf{Processing Pipeline:}
\begin{enumerate}
    \item Capture aligned depth + color frames from RealSense at 30 FPS
    \item Retrieve RealSense color camera intrinsics $\mathbf{K}_{\text{RS}}$
    \item Fetch AVP camera intrinsics $\mathbf{K}_{\text{AVP}}$ from main API
    \item Load or derive extrinsics $(\mathbf{R}_{\text{AVP←RS}}, \mathbf{t}_{\text{AVP←RS}})$
    \item Reproject depth via 3D transformation and Z-buffering (detailed in §\ref{sec:realsense_transform})
    \item Normalize aligned depth to 0-255 disparity map
\end{enumerate}

\textbf{Performance:} 33ms hardware capture + 20-50ms reprojection ≈ 50-80ms total

\subsubsection{RealSense to AVP Coordinate Transformation}
\label{sec:realsense_transform}

The RealSense depth sensor and the Apple Vision Pro main camera occupy different physical positions and orientations. To align RealSense depth with AVP's RGB view, we perform a full 3D reprojection with Z-buffering.

\paragraph{Mathematical Formulation}

Let:
\begin{itemize}
    \item $\mathbf{K}_{\text{RS}} \in \mathbb{R}^{3 \times 3}$ -- RealSense color camera intrinsics
    \item $\mathbf{K}_{\text{AVP}} \in \mathbb{R}^{3 \times 3}$ -- AVP camera intrinsics
    \item $\mathbf{R}_{\text{AVP←RS}} \in \text{SO}(3)$ -- Rotation from RS color frame to AVP frame
    \item $\mathbf{t}_{\text{AVP←RS}} \in \mathbb{R}^3$ -- Translation from RS color frame to AVP frame
    \item $D_{\text{RS}}(u, v)$ -- Depth map in RealSense color frame (meters)
\end{itemize}

\textbf{Step 1: Back-projection to 3D in RealSense Frame}

For each pixel $(u, v)$ with depth $Z = D_{\text{RS}}(u, v)$:
\[
\begin{bmatrix} X \\ Y \\ Z \end{bmatrix}_{\text{RS}}
= Z \cdot \mathbf{K}_{\text{RS}}^{-1} \begin{bmatrix} u \\ v \\ 1 \end{bmatrix}
= \begin{bmatrix}
\frac{(u - c_x) \cdot Z}{f_x} \\
\frac{(v - c_y) \cdot Z}{f_y} \\
Z
\end{bmatrix}_{\text{RS}}
\]

where $f_x, f_y, c_x, c_y$ are intrinsic parameters from $\mathbf{K}_{\text{RS}}$.

\textbf{Step 2: Transform to AVP Frame}

Apply rigid transformation:
\[
\begin{bmatrix} X \\ Y \\ Z \end{bmatrix}_{\text{AVP}}
= \mathbf{R}_{\text{AVP←RS}} \begin{bmatrix} X \\ Y \\ Z \end{bmatrix}_{\text{RS}}
+ \mathbf{t}_{\text{AVP←RS}}
\]

\textbf{Step 3: Project to AVP Image Plane}

For valid points with $Z_{\text{AVP}} > \epsilon$:
\[
\begin{bmatrix} u' \\ v' \\ 1 \end{bmatrix}
= \frac{1}{Z_{\text{AVP}}} \mathbf{K}_{\text{AVP}} \begin{bmatrix} X \\ Y \\ Z \end{bmatrix}_{\text{AVP}}
\]

\[
u' = \frac{f'_x \cdot X_{\text{AVP}}}{Z_{\text{AVP}}} + c'_x, \quad
v' = \frac{f'_y \cdot Y_{\text{AVP}}}{Z_{\text{AVP}}} + c'_y
\]

\textbf{Step 4: Z-Buffer Splatting}

Multiple RealSense pixels may project to the same AVP pixel. To resolve occlusion correctly, we maintain a depth buffer:

\begin{algorithmic}
\STATE Initialize: $D_{\text{AVP}}(u', v') \leftarrow \infty$ for all pixels
\FOR{each valid 3D point with projection $(u'_i, v'_i, Z_i)$}
    \IF{$Z_i < D_{\text{AVP}}(u'_i, v'_i)$}
        \STATE $D_{\text{AVP}}(u'_i, v'_i) \leftarrow Z_i$ \quad \textcolor{gray}{\textit{(keep nearest)}}
    \ENDIF
\ENDFOR
\STATE Replace $\infty$ with 0 for unfilled pixels
\end{algorithmic}

This ensures proper occlusion handling where nearer surfaces occlude farther ones.

\paragraph{Extrinsics Calibration}

The transformation $(\mathbf{R}_{\text{AVP←RS}}, \mathbf{t}_{\text{AVP←RS}})$ can be obtained via:

\textbf{Option 1: Manual Calibration}

Pre-calibrate using a checkerboard or ArUco board visible to both cameras simultaneously. Store in \texttt{extrinsics\_avp\_from\_rs\_color.json}:

\begin{verbatim}
{
  "R": [[r11, r12, r13], [r21, r22, r23], [r31, r32, r33]],
  "t": [tx, ty, tz]
}
\end{verbatim}

\textbf{Option 2: Runtime Estimation}

If both cameras detect the same ArUco board simultaneously:
\begin{enumerate}
    \item Solve PnP to get $\mathbf{T}_{\text{RS}}^{\text{board}}$ and $\mathbf{T}_{\text{AVP}}^{\text{board}}$
    \item Compute relative transform: $\mathbf{T}_{\text{AVP←RS}} = \mathbf{T}_{\text{AVP}}^{\text{board}} \cdot (\mathbf{T}_{\text{RS}}^{\text{board}})^{-1}$
    \item Extract $\mathbf{R}_{\text{AVP←RS}}$ and $\mathbf{t}_{\text{AVP←RS}}$
\end{enumerate}

This method is less stable but enables automatic online calibration.

\subsubsection{ArUco Board Configuration}

The system uses a 3×4 grid of ArUco markers (12 total) from dictionary DICT\_4X4\_50:

\begin{itemize}
    \item \textbf{Marker size:} 30mm × 30mm
    \item \textbf{Separation:} 10mm
    \item \textbf{Total board size:} 190mm × 150mm
\end{itemize}

\paragraph{3D Marker Layout}

Marker ID $i \in \{0, 1, \ldots, 11\}$ maps to board position:

\[
\text{row} = \lfloor i / 4 \rfloor, \quad \text{col} = i \mod 4
\]

Corner positions in board frame (Z=0 plane):

\[
\mathbf{p}_{\text{corner}_k}^{(i)} = \begin{bmatrix}
x_0 + k_x \cdot s \\
y_0 + k_y \cdot s \\
0
\end{bmatrix}, \quad k \in \{0, 1, 2, 3\}
\]

where:
\[
x_0 = \text{col} \cdot (s + d), \quad
y_0 = \text{row} \cdot (s + d), \quad
s = 0.030\,\text{m}, \quad
d = 0.010\,\text{m}
\]

and $(k_x, k_y) \in \{(0,0), (1,0), (1,1), (0,1)\}$ for corners in order.

\paragraph{Pose Estimation via PnP}

Given detected marker corners in image $\{(u_j, v_j)\}$ and corresponding 3D points $\{\mathbf{p}_j\}$:

\begin{enumerate}
    \item Concatenate all 2D and 3D correspondences from detected markers
    \item Solve PnP with IPPE algorithm: $\min_{\mathbf{R}, \mathbf{t}} \sum_j \| \pi(\mathbf{R}\mathbf{p}_j + \mathbf{t}) - (u_j, v_j) \|^2$
    \item Fallback to standard PnP if IPPE fails
    \item Return rotation vector $\mathbf{r} \in \mathbb{R}^3$ and translation $\mathbf{t} \in \mathbb{R}^3$
\end{enumerate}

\subsubsection{Main API Server}

The \texttt{main\_api.py} Flask server provides RESTful endpoints for the entire system:

\paragraph{Frame Processing Endpoints}

\begin{itemize}
    \item \texttt{POST /receive\_frame} -- Accept RGB frame, process through CV pipeline
    \item \texttt{GET /rgb\_frame} -- Retrieve last raw frame
    \item \texttt{GET /detected\_frame} -- Retrieve frame with ArUco markers annotated
    \item \texttt{GET /intrinsics} -- Camera intrinsics matrix
    \item \texttt{GET /pose} -- Last detected board pose (rvec, tvec)
    \item \texttt{GET /mask} -- ROI mask as base64 PNG
    \item \texttt{GET /depth} -- Depth map as base64 PNG
    \item \texttt{GET /disparity} -- Disparity map as base64 PNG
\end{itemize}

\paragraph{Head Pose Tracking}

\begin{itemize}
    \item \texttt{POST /head\_pose} -- Receive head tracking data from Vision Pro
    \item \texttt{GET /head\_pose} -- Retrieve last head pose with age
\end{itemize}

Head pose format (JSON):
\begin{verbatim}
{
  "position": [x, y, z],           // meters
  "rotation": [roll, pitch, yaw],  // radians
  "quaternion": [qx, qy, qz, qw],  // unit quaternion
  "timestamp": 1699876543.21,      // Unix time
  "confidence": 0.95
}
\end{verbatim}

\paragraph{Model Management}

\begin{itemize}
    \item \texttt{GET /models} -- List available .ply CAD models
    \item \texttt{GET /model?name=X} -- Download specific model (base64)
    \item \texttt{POST /select\_model} -- Set active model for pose estimation
\end{itemize}

\paragraph{AVP Pose Estimation Endpoint}

\texttt{POST /avp\_pose} -- Primary endpoint for Vision Pro pose requests

\textbf{Request payload:}
\begin{verbatim}
{
  "rgb_frame": "data:image/jpeg;base64,...",
  "camera_matrix": [[fx,0,cx],[0,fy,cy],[0,0,1]],
  "depth_map": "data:image/png;base64,...",  // optional
  "mask": "data:image/jpeg;base64,...",      // optional
  "model_name": "cube.ply"                   // optional
}
\end{verbatim}

\textbf{Processing flow:}
\begin{enumerate}
    \item Validate request and load selected CAD model
    \item Use provided depth/mask or fetch from CV pipeline cache
    \item Forward to pose estimation API with full payload
    \item Receive 4×4 transformation matrix $\mathbf{T}_{\text{object}}^{\text{camera}}$
    \item Apply head pose correction (if available)
    \item Return corrected pose to client
\end{enumerate}

\subsubsection{Head Pose Correction}

To compensate for head motion between frame capture and pose rendering, the backend applies a correction transform.

\paragraph{Quaternion to Rotation Matrix}

Vision Pro provides head orientation as quaternion $\mathbf{q} = (q_x, q_y, q_z, q_w)$. Convert to rotation matrix:

\[
\mathbf{R}_{\text{head}} = \begin{bmatrix}
1 - 2(q_y^2 + q_z^2) & 2(q_x q_y - q_w q_z) & 2(q_x q_z + q_w q_y) \\
2(q_x q_y + q_w q_z) & 1 - 2(q_x^2 + q_z^2) & 2(q_y q_z - q_w q_x) \\
2(q_x q_z - q_w q_y) & 2(q_y q_z + q_w q_x) & 1 - 2(q_x^2 + q_y^2)
\end{bmatrix}
\]

\paragraph{Transformation Composition}

Build head transformation matrix:
\[
\mathbf{T}_{\text{head}} = \begin{bmatrix}
\mathbf{R}_{\text{head}} & \mathbf{t}_{\text{head}} \\
\mathbf{0}^T & 1
\end{bmatrix}
\]

Apply correction:
\[
\mathbf{T}_{\text{corrected}} = \mathbf{T}_{\text{head}} \cdot \mathbf{T}_{\text{object}}
\]

This compensates for head motion, ensuring the overlay remains anchored to the physical object even as the user moves.

\textbf{Staleness Check:} Head pose data older than 1 second is ignored to prevent incorrect corrections.

\subsection{VisionOS Client Application}

\subsubsection{Application Structure}

The \texttt{PoseOverlayApp} is built with SwiftUI and RealityKit, leveraging Vision Pro's spatial computing capabilities.

\textbf{Core Components:}
\begin{itemize}
    \item \texttt{AppModel} -- Global state (base URL, selected model, settings)
    \item \texttt{ImmersiveSpaceView} -- AR scene with RealityKit rendering
    \item \texttt{PoseService} -- Networking layer for backend communication
    \item \texttt{SensorDataModel} -- ARKit integration for head tracking
    \item \texttt{LogStore} -- Timestamped event logging for debugging
\end{itemize}

\subsubsection{Pose Service and Snapshot Retrieval}

The \texttt{PoseService} enum provides a stateless API for fetching pose data from the Python backend.

\paragraph{Snapshot Acquisition}

Before requesting a pose, the client fetches a complete snapshot from the pipeline:

\begin{verbatim}
struct PipelineSnapshot {
    let rgbFrame: String          // base64 JPEG
    let cameraMatrix: [[Double]]  // 3x3 intrinsics
    let mask: String?             // base64 from /mask
    let depthMap: String?         // base64 PNG from /disparity
}
\end{verbatim}

All endpoints are queried concurrently using Swift's \texttt{async/await}:

\begin{verbatim}
async let rgbTask = fetchRGBFrame(baseURL: baseURL)
async let intrinsicsTask = fetchIntrinsics(baseURL: baseURL)
async let maskTask = fetchOptionalMask(baseURL: baseURL)
async let depthTask = fetchOptionalDisparity(baseURL: baseURL)
\end{verbatim}

\paragraph{Pose Request}

Payload is assembled and sent to \texttt{/avp\_pose}:

\begin{verbatim}
POST /avp_pose
{
  "rgb_frame": snapshot.rgbFrame,
  "camera_matrix": snapshot.cameraMatrix,
  "depth_map": snapshot.depthMap,
  "mask": snapshot.mask,
  "model_name": selectedModel
}
\end{verbatim}

Response contains transformation matrices as 4×4 nested arrays.

\subsubsection{Coordinate System Transformations}

\paragraph{OpenCV to RealityKit Conversion}

The backend returns poses in OpenCV convention (right-handed, Z-forward, Y-down). RealityKit uses a different convention (right-handed, Z-backward, Y-up). The client applies a conversion:

\[
\mathbf{T}_{\text{RK}} = \mathbf{C} \cdot \mathbf{T}_{\text{CV}}
\]

where the conversion matrix is:

\[
\mathbf{C} = \begin{bmatrix}
1 & 0 & 0 & 0 \\
0 & -1 & 0 & 0 \\
0 & 0 & -1 & 0 \\
0 & 0 & 0 & 1
\end{bmatrix}
\]

This flips the Y and Z axes to match RealityKit's coordinate frame.

\paragraph{SE(3) Representation}

All transformations are members of the special Euclidean group $\text{SE}(3)$:

\[
\mathbf{T} = \begin{bmatrix}
\mathbf{R} & \mathbf{t} \\
\mathbf{0}^T & 1
\end{bmatrix} \in \text{SE}(3), \quad
\mathbf{R} \in \text{SO}(3), \quad
\mathbf{t} \in \mathbb{R}^3
\]

Composition: $\mathbf{T}_1 \cdot \mathbf{T}_2$ represents applying $\mathbf{T}_2$ then $\mathbf{T}_1$

Inverse: $\mathbf{T}^{-1} = \begin{bmatrix} \mathbf{R}^T & -\mathbf{R}^T \mathbf{t} \\ \mathbf{0}^T & 1 \end{bmatrix}$

\subsubsection{Immersive Space and Rendering}

\paragraph{Scene Hierarchy}

\begin{verbatim}
HeadAnchor (AnchorEntity .head)
├── PoseOverlayRoot (Entity at [0, 0, -1])
│   └── poseArrow_0..N (ModelEntity per pose)
└── HeadAxes (Visualization entity)
    ├── X-axis (red)
    ├── Y-axis (green)
    └── Z-axis (blue)
\end{verbatim}

The root container is positioned 1 meter in front of the user's head and remains head-locked. Pose entities are added as children with transforms from the backend.

\paragraph{Polling Loop}

The app uses a background task to continuously poll for pose updates:

\begin{verbatim}
while !Task.isCancelled {
    let transforms = try await PoseService.fetchTransforms(
        baseURL: baseURL,
        modelName: selectedModel
    )
    await MainActor.run {
        update(container: container, with: transforms)
    }
    try? await Task.sleep(nanoseconds: 2_000_000_000) // 2s
}
\end{verbatim}

Each iteration:
\begin{enumerate}
    \item Fetches snapshot from backend
    \item Requests pose estimation
    \item Converts matrices to RealityKit format
    \item Updates scene entities on main actor
\end{enumerate}

\paragraph{Model Rendering}

For each transformation matrix $\mathbf{T}_i$:
\begin{enumerate}
    \item Create or retrieve arrow entity (visual placeholder for CAD model)
    \item Convert matrix to \texttt{simd\_float4x4}
    \item Apply RealityKit coordinate conversion
    \item Set entity transform: \texttt{entity.transform = Transform(matrix: T)}
    \item Add to scene container
\end{enumerate}

\subsubsection{Sensor Integration}

\paragraph{ARKit Head Tracking}

The \texttt{SensorDataModel} wraps ARKit's world tracking:

\begin{itemize}
    \item Provides real-time head position and orientation
    \item Updates UI-bound \texttt{@Published} properties
    \item Sends periodic updates to backend \texttt{/head\_pose} endpoint
\end{itemize}

Head orientation is represented as a quaternion and synchronized with the axes visualization entity for debugging.

\subsection{Integration and Workflow}

\subsubsection{End-to-End Dataflow}

\textbf{Phase 1: Initialization}
\begin{enumerate}
    \item Vision Pro app starts, user enters backend URL
    \item App queries \texttt{/models} to populate model picker
    \item User selects CAD model, app sends \texttt{POST /select\_model}
    \item Immersive space opens, polling task starts
\end{enumerate}

\textbf{Phase 2: Runtime Loop}
\begin{enumerate}
    \item Backend CV pipeline processes screen-captured frames (or RealSense if available)
    \item Pipeline caches: RGB, intrinsics, mask, depth/disparity, ArUco pose
    \item Vision Pro polls backend every 2s
    \item \texttt{PoseService} fetches snapshot (RGB, K, mask, depth) concurrently
    \item Snapshot + model sent to \texttt{/avp\_pose}
    \item Backend forwards to pose API, receives $\mathbf{T}_{\text{object}}$
    \item Backend applies head pose correction if available
    \item Client receives corrected pose, converts to RealityKit coordinates
    \item RealityKit renders CAD overlay at pose location
\end{enumerate}

\textbf{Phase 3: Head Tracking Correction}
\begin{enumerate}
    \item Vision Pro continuously tracks head pose via ARKit
    \item \texttt{SensorDataModel} sends updates to \texttt{POST /head\_pose}
    \item Backend caches latest head pose with timestamp
    \item When pose estimate arrives, backend checks head pose freshness
    \item If $\Delta t < 1\text{s}$, apply correction: $\mathbf{T}_{\text{out}} = \mathbf{T}_{\text{head}} \cdot \mathbf{T}_{\text{object}}$
    \item Corrected pose sent to client
\end{enumerate}

\subsubsection{Concurrency Design}

\paragraph{Backend: Threaded Processing}

\begin{itemize}
    \item \textbf{Main thread:} Flask request handlers
    \item \textbf{Pipeline state:} Protected by threading locks
    \item \textbf{Frame capture:} Independent thread (screen\_capture.py or RealSense)
    \item \textbf{Depth inference:} Asynchronous GPU execution
\end{itemize}

\paragraph{Client: Swift Concurrency}

\begin{itemize}
    \item \textbf{MainActor:} UI updates, RealityKit scene mutations
    \item \textbf{Background tasks:} Polling loop, network requests
    \item \textbf{async/await:} Concurrent snapshot fetching
    \item \textbf{Actors:} \texttt{SensorDataModel} for ARKit isolation
\end{itemize}

\subsection{Performance Analysis}

\subsubsection{Latency Breakdown}

\begin{table}[H]
\centering
\begin{tabular}{|l|r|r|}
\hline
\textbf{Stage} & \textbf{Transformers Mode} & \textbf{RealSense Mode} \\
\hline
Frame capture & 33ms & 33ms \\
ArUco detection & 5-10ms & 5-10ms \\
Depth acquisition & 50-100ms (GPU) & 33ms (hw) \\
Depth alignment & N/A & 20-50ms \\
Mask extraction & 2-5ms & 2-5ms \\
PNG encoding & 5-10ms & 5-10ms \\
Network (AVP↔backend) & 10-50ms & 10-50ms \\
Pose API inference & 100-500ms & 100-500ms \\
Head pose correction & <1ms & <1ms \\
RealityKit rendering & 16ms (60fps) & 16ms (60fps) \\
\hline
\textbf{Total (approx)} & \textbf{220-730ms} & \textbf{190-660ms} \\
\hline
\end{tabular}
\caption{Per-frame latency breakdown for both depth modes}
\label{tab:latency}
\end{table}

\subsubsection{Optimization Strategies}

\begin{enumerate}
    \item \textbf{Depth caching:} Reuse disparity across multiple pose requests if intrinsics unchanged
    \item \textbf{ROI cropping:} Reduce inference resolution to speed up depth estimation
    \item \textbf{Model quantization:} Use INT8 quantized depth models for faster inference
    \item \textbf{Parallel requests:} Client fetches snapshot components concurrently
    \item \textbf{GPU batching:} Process multiple frames in batch for higher throughput
\end{enumerate}

\subsection{Configuration and Deployment}

\subsubsection{Backend Configuration}

Edit \texttt{app\_config.py}:

\begin{verbatim}
APP_CONFIG = {
    "main_api": {"host": "0.0.0.0", "port": 5000},
    "pose_api": {"base_url": "http://pose-server:9000"},
    "defaults": {
        "use_realsense": False,  # Toggle depth mode
        "model_name": "cube.ply",
        "use_random_pose": False  # Use real pose API
    }
}
\end{verbatim}

\subsubsection{Launch Sequence}

\textbf{Terminal 1 -- Backend API:}
\begin{verbatim}
cd AW19
python main_api.py
# Listens on http://0.0.0.0:5000
\end{verbatim}

\textbf{Terminal 2 -- Frame Capture (if not using RealSense):}
\begin{verbatim}
python screen_capture.py
# Sends frames to /receive_frame
\end{verbatim}

\textbf{Terminal 3 -- Debug Viewer (optional):}
\begin{verbatim}
python full_python_pipeline/tk_debugging_unified.py
# Shows: RGB, Mask, Depth, Disparity, Pose
\end{verbatim}

\textbf{Vision Pro:}
\begin{enumerate}
    \item Build and run \texttt{PoseOverlayApp} in Xcode
    \item Enter backend URL (e.g., \texttt{http://192.168.1.100:5000})
    \item Select model from picker
    \item Tap "Open Immersive Space"
    \item Point device at ArUco board and wait for pose overlay
\end{enumerate}

\subsection{Summary}

This chapter presented the complete implementation spanning:

\begin{itemize}
    \item \textbf{Python Backend:} Unified CV pipeline with dual depth support, ArUco detection, head pose correction
    \item \textbf{RealSense Integration:} 3D reprojection with Z-buffering for metric depth alignment
    \item \textbf{VisionOS Client:} RealityKit-based AR rendering with snapshot fetching and coordinate conversion
    \item \textbf{Geometric Transformations:} SE(3) pose representations, quaternion conversions, frame transformations
    \item \textbf{System Integration:} RESTful API design, concurrent processing, latency optimization
\end{itemize}

The modular architecture enables flexible deployment across different hardware configurations while maintaining a consistent interface for pose estimation and AR visualization.
